%======================================================================
\section{Introduction}
\label{sec:intro}

The most direct way to ask what a component of a network contributes to an output is to remove the component and measure the change. A first-order estimate replaces each removal with a linear prediction computed from the derivative of the output at the intact network. The idea goes back to the relevance and saliency estimates used to prune networks~\citep{mozer1989skeletonization,lecun1990optimal,molchanov2017pruning,molchanov2019importance}, and in interpretability research it underlies attribution patching, the first-order approximation of activation patching in language models~\citep{nanda2023attribution,syed2024attribution,kramar2024atp,hanna2024faith}. Such estimates serve two purposes. They can screen, picking out the components whose removal matters, and they can estimate, giving the size and sign of each effect. The two can come apart: a linear prediction can rank components well and still misjudge how much each one changes the output, particularly when the removal is large compared with the range over which the network responds linearly~\citep{nanda2023attribution,zhang2026attribution}.

We separate the two purposes, and then look for the cause of the failure, in a setting where every removal can be simulated exactly and the components have biological identities. The FlyVis models of \citet{lappalainen2024connectome} are \NModels{} recurrent networks of the fly visual system whose connectivity comes from the connectome. Each model has \NetNodes{} neurons of \NetTypes{} types and was trained on optic-flow estimation; the models share the connectome and differ in their trained parameters. Silencing a cell type is a well-defined intervention whose effects can travel along indirect and recurrent paths, and the motion detectors T4 and T5 have known direction selectivity and well-studied inputs~\citep{borst2020fly,strother2017emergence,ammer2015functional,serbe2016comprehensive}. The FlyVis authors themselves silenced the four main inputs to T4c in their models and compared the outcome qualitatively with experiments~\citep{lappalainen2024review}.

\textbf{Study 1} silenced each cell type at full and at half strength in every model and measured the change in the direction-selectivity index (DSI) of the eight T4 and T5 subtypes, using ON and OFF edges moving in \NDirections{} directions at two speeds. It compared these exact changes with first-order estimates and, for screening, with rankings computed from the connectome, from the trained weights and from the input currents of the intact model. \textbf{Study 2} asked why the estimates misjudge the effects. Three explanations were open:
\begin{itemize}
  \item \textbf{The baseline.} Silencing drives a cell type's output to zero, far from where the network operates. Replacing each cell's output by its own average removes the cell's stimulus-dependent signal but keeps its average drive, a smaller perturbation. The choice of replacement value is known to change what patching experiments find~\citep{heimersheim2024patching,miller2024metrics}, and in the same FlyVis models an estimate whose baseline differs from the intervention loses the size of the effects~\citep{dhawan2026flyhack}.
  \item \textbf{The readout.} The direction-selectivity index takes the temporal peak of the rectified voltage in each trajectory and then divides a vector length by a sum of responses. Both steps are nonlinear, and the peak is not smooth.
  \item \textbf{Curvature.} The first-order estimate omits the second-order term of the Taylor expansion in the silencing strength, which \citet{zhang2026attribution} add to attribution patching in language models.
\end{itemize}
Study 2 tested each explanation by changing one thing at a time from Study 1's setting, on the same models, stimuli, readouts and cell types, and it removed one source of nonlinearity at a time from Study 1's estimate. Each study fixed its analysis on four development models before the other \NConfModels{} models, the confirmation set, were analysed, and Study 2 hashed \PredPrimaryTotal{} predictions before any of its runs. Our findings are:
\begin{itemize}
  \item \textbf{The first-order estimate is a good screen.} At full silencing, a mean fraction of \PrecFOFull{} of the ten cell types it ranked highest had large effects ($|\Delta\DSI|\ge\LargeThreshold$), against \PrecAnatFull{} for anatomical path weights through three synaptic steps, the prespecified comparator, and \PrecAnatFourFull{} for the best anatomical variant. Counting abolished responses as large effects and scoring every responsive readout condition, the fractions were \InclFOFull{} and \InclAnatFull{} (\cref{sec:screening}).
  \item \textbf{Its effect sizes are worse than a prediction of zero change.} At full silencing its mean absolute error was \MaeFOFull{}, against \MaeZeroFull{} for predicting zero change, and it was the larger error in \MaeWorseFull{} of \NConfModels{} models; \WrongSignFull\% of large changes had the wrong sign. Among large effects, its ranking was no better than anatomy (Spearman correlation \RhoLargeFOFull{} against \RhoLargeAnatFull{}; \cref{sec:effects,sec:ranking}).
  \item \textbf{The readout's nonlinearity, not the network's, accounts for the error.} With the network linearised and the index computed exactly, the error ratio (the estimate's mean absolute error divided by that of predicting zero change) falls from \RatioRefFull{} to \RatioCfReadFull{} \RatioCICfReadFull{}, below 1 in \BelowOneCfRead{} of \NConfModels{} models, and wrong signs fall to \WrongCfReadFull\%. Linearising the readout on the exact trajectories gives \RatioCfNetFull{} (\cref{sec:mechanism}).
  \item \textbf{Within the readout, the ratio matters most.} Applying the exact ratio to linearised responses (floored at zero) already beats predicting zero change (\RatioSplitRatioFull{}); making the responses exact while linearising the ratio makes the estimate far worse (\RatioSplitRespFull{}; \cref{sec:split}).
  \item \textbf{Of the three single changes, only mean replacement improves on the first-order estimate} (error ratio \RatioAFull{}, change \ChangeAFull{} \ChangeCIAFull{}), but at full strength it is still not distinguishable from predicting zero change. The mean-response index changes nothing detectable (\RatioBFull{}), and the second-order correction makes the estimates far worse (\RatioCFull{}); a post hoc analysis traces this to the ratio, whose Taylor series diverges where silencing changes the responses by a large fraction (\cref{sec:arms,sec:second}).
  \item \textbf{The estimates cost about as much as exact silencing at one strength.} Silencing the \SimulatedTypes{} cell types with a path to the readouts took \TimeExactFull{}~s per model at one strength and \TimeExactBoth{}~s at both, and the derivatives for the eight readouts took \TimeGradient{}~s (\cref{sec:cost}).
\end{itemize}
The contribution is an evaluation of a known approximation, and a diagnosis of its failure, with exact silencing in the model as the reference. Biological experiments provide context for that reference, not a test of it.
